\documentclass[a4paper,12pt]{article}

\usepackage[T2A]{fontenc}
\usepackage[utf8]{inputenc}
\usepackage[russian]{babel}

\usepackage{amsmath, amssymb, amsthm, mathtools}
\usepackage{helvet}
\renewcommand{\familydefault}{\sfdefault}
\usepackage{icomma}
\usepackage{geometry}
\usepackage{microtype}
\usepackage{tikz}
\usetikzlibrary{arrows.meta,calc}
\usepackage{tkz-euclide}
\usepackage{pgfplots}
\pgfplotsset{compat=1.18}
\usepackage{tabularx, booktabs, longtable}
\usepackage{enumitem}
\usepackage{hyperref}
\usepackage{parskip}
\usepackage{fancyhdr}
\usepackage{setspace}
\usepackage{xcolor}
\usepackage{titlesec}

\geometry{left=18mm,right=18mm,top=18mm,bottom=20mm}
\hypersetup{hidelinks}
\onehalfspacing

\definecolor{accent}{HTML}{008080}
\definecolor{darkaccent}{HTML}{004D4D}
\definecolor{textgray}{HTML}{333333}
\definecolor{softgray}{HTML}{F4F6F6}

\color{textgray}

\pagestyle{fancy}
\fancyhf{}
\fancyhead[L]{\small Тригонометрические уравнения}
\fancyhead[R]{\small Тимофей Васильченко}
\fancyfoot[C]{\small \thepage}
\renewcommand{\headrulewidth}{0.4pt}
\renewcommand{\headrule}{\hbox to\headwidth{\color{accent}\leaders\hrule height \headrulewidth\hfill}}

\titleformat{\section}
  {\Large\bfseries\color{darkaccent}}
  {\thesection}{0.8em}{}

\titleformat{\subsection}
  {\large\bfseries\color{accent}}
  {\thesubsection}{0.8em}{}

\setlist[itemize]{leftmargin=1.4em,itemsep=0.25em,topsep=0.25em}
\setlist[enumerate]{leftmargin=1.7em,itemsep=0.35em,topsep=0.35em}

\newcolumntype{Y}{>{\raggedright\arraybackslash}X}

\newcommand{\R}{\mathbb{R}}
\newcommand{\Z}{\mathbb{Z}}
\newcommand{\ans}[1]{\(\displaystyle #1\)}
\newcommand{\arcctg}{\operatorname{arcctg}}

\begin{document}

\begin{titlepage}
\centering

\vspace*{2.1cm}

{\color{accent}\rule{0.75\linewidth}{0.8pt}\par}
\vspace{1.2cm}

{\Huge\bfseries\color{darkaccent} Чек-лист\par}
\vspace{0.5cm}

{\Large\bfseries Тригонометрические уравнения\par}
\vspace{0.25cm}
{\large Корни, равносильные переходы, модули и отбор решений\par}

\vspace{1.8cm}

{\large ЕГЭ\par}

\vfill

{\large Лёвин Артём Александрович\\[0.2cm]
эксклюзивно для Тимофея Васильченко\par}

\vspace{0.8cm}

{\large \today\par}

\vspace*{1.5cm}

\begin{tikzpicture}[remember picture,overlay]
  \draw[accent,line width=1.1pt]
    ([xshift=2.0cm,yshift=2.1cm]current page.south west)
    .. controls ([xshift=5.2cm,yshift=3.0cm]current page.south west)
    and ([xshift=-5.2cm,yshift=1.2cm]current page.south east)
    .. ([xshift=-2.0cm,yshift=2.3cm]current page.south east);
  \draw[darkaccent,line width=0.55pt]
    ([xshift=2.8cm,yshift=1.55cm]current page.south west)
    .. controls ([xshift=6.5cm,yshift=2.0cm]current page.south west)
    and ([xshift=-6.0cm,yshift=3.1cm]current page.south east)
    .. ([xshift=-2.8cm,yshift=1.65cm]current page.south east);
\end{tikzpicture}

\end{titlepage}

\section{Главная идея}

В уравнениях с корнями, модулями и тригонометрическими функциями нельзя просто механически возводить обе части в квадрат или раскрывать модуль. Каждый такой переход должен быть равносильным.

Главная задача решения:

\[
\text{не потерять корни и не добавить лишние.}
\]

Для этого нужно:

\begin{itemize}
  \item записывать ограничения перед возведением в квадрат;
  \item раскрывать модуль по определению;
  \item проверять найденные значения тригонометрических функций;
  \item отбирать решения на тригонометрической окружности;
  \item аккуратно работать с отрицательными значениями в формулах для \(\cos x\), \(\sin x\), \(\tg x\).
\end{itemize}

\section{Уравнение вида \(\sqrt{f(x)}=g(x)\)}

\subsection{Равносильный переход}

Если слева корень, а справа выражение без корня, то используем переход:

\[
\sqrt{f(x)}=g(x)
\quad\Longleftrightarrow\quad
\begin{cases}
g(x)\ge 0,\\
f(x)=g^2(x).
\end{cases}
\]

Ограничение ставится именно на правую часть \(g(x)\), потому что корень чётной степени не может быть отрицательным.

Условие \(f(x)\ge0\) отдельно писать не требуется: из равенства

\[
f(x)=g^2(x)
\]

автоматически следует, что \(f(x)\ge0\).

\subsection{Пример}

Решим уравнение:

\[
\sqrt{\cos^2 x}=\frac12-\cos x.
\]

Так как \(\sqrt{\cos^2 x}=|\cos x|\), можно решить через модуль. Но для демонстрации равносильного перехода рассмотрим именно схему с корнем:

\[
\sqrt{\cos^2 x}=\frac12-\cos x
\quad\Longleftrightarrow\quad
\begin{cases}
\frac12-\cos x\ge0,\\
\cos^2x=\left(\frac12-\cos x\right)^2.
\end{cases}
\]

Первое условие:

\[
\cos x\le\frac12.
\]

Второе условие:

\[
\cos^2x=\frac14-\cos x+\cos^2x.
\]

Сокращаем \(\cos^2x\):

\[
0=\frac14-\cos x,
\]

\[
\cos x=\frac14.
\]

Теперь проверяем ограничение:

\[
\frac14\le\frac12.
\]

Ограничение выполнено, значит:

\[
\cos x=\frac14.
\]

Ответ:

\[
x=\pm\arccos\frac14+2\pi n,\qquad n\in\Z.
\]

\section{Когда нужно сравнивать полученное число}

Иногда после замены получается значение, которое нужно сравнить с ограничением.

Например, если получено

\[
\cos x=2-\sqrt2,
\]

а ограничение требует

\[
\cos x\le\frac12,
\]

то нужно проверить:

\[
2-\sqrt2\stackrel{?}{\le}\frac12.
\]

Сравним:

\[
2-\sqrt2\le\frac12
\quad\Longleftrightarrow\quad
-\sqrt2\le-\frac32.
\]

Домножаем на \(-1\), знак меняется:

\[
\sqrt2\ge\frac32.
\]

Возводим обе части в квадрат:

\[
2\ge\frac94.
\]

Это неверно, значит

\[
2-\sqrt2>\frac12.
\]

Такое значение не подходит под ограничение \(\cos x\le\frac12\).

\section{Уравнение вида \(\sqrt{f(x)}=\sqrt{g(x)}\)}

\subsection{Равносильный переход}

Если корни стоят с обеих сторон, то достаточно записать неотрицательность одной части и приравнять подкоренные выражения:

\[
\sqrt{f(x)}=\sqrt{g(x)}
\quad\Longleftrightarrow\quad
\begin{cases}
f(x)\ge0,\\
f(x)=g(x).
\end{cases}
\]

Можно также писать:

\[
\sqrt{f(x)}=\sqrt{g(x)}
\quad\Longleftrightarrow\quad
\begin{cases}
g(x)\ge0,\\
f(x)=g(x).
\end{cases}
\]

Обычно выбирают то ограничение, которое проще отбирать.

\subsection{Пример}

Решим уравнение:

\[
\sqrt{\sin x}=\sqrt{\cos x}.
\]

Переходим к системе:

\[
\begin{cases}
\sin x\ge0,\\
\sin x=\cos x.
\end{cases}
\]

Из уравнения:

\[
\sin x=\cos x.
\]

Так как \(\cos x\ne0\) для решений этого уравнения, делим на \(\cos x\):

\[
\tg x=1.
\]

Отсюда:

\[
x=\frac{\pi}{4}+\pi n,\qquad n\in\Z.
\]

Теперь отбираем по условию \(\sin x\ge0\). На окружности подходят точки в верхней полуплоскости.

\begin{center}
\begin{tikzpicture}[scale=1.15,>=Stealth]
  \draw[->] (-1.35,0)--(1.45,0) node[right] {\(\cos x\)};
  \draw[->] (0,-1.35)--(0,1.45) node[above] {\(\sin x\)};
  \draw[darkaccent,line width=0.9pt] (0,0) circle (1);
  \draw[accent,line width=2pt] (1,0) arc (0:180:1);

  \coordinate (A) at ({cos(45)},{sin(45)});
  \coordinate (B) at ({cos(225)},{sin(225)});

  \fill[accent] (A) circle (2pt);
  \fill[textgray] (B) circle (1.8pt);

  \draw[accent,dashed] (0,0)--(A);
  \draw[textgray,dashed] (0,0)--(B);

  \node[above right] at (A) {\(\frac{\pi}{4}\)};
  \node[below left] at (B) {\(\frac{5\pi}{4}\)};
  \node[above left] at (-0.95,0.1) {\small \(\sin x\ge0\)};
\end{tikzpicture}
\end{center}

Ответ:

\[
x=\frac{\pi}{4}+2\pi n,\qquad n\in\Z.
\]

\section{Как решать уравнение после возведения в квадрат}

После равносильного перехода часто появляется квадратное уравнение относительно \(\sin x\) или \(\cos x\).

\subsection{Замена}

Если уравнение содержит только \(\cos x\), удобно сделать замену:

\[
t=\cos x,\qquad -1\le t\le1.
\]

Если уравнение содержит только \(\sin x\), удобно сделать замену:

\[
t=\sin x,\qquad -1\le t\le1.
\]

Ограничение \(-1\le t\le1\) обязательно помнить, потому что синус и косинус не могут быть больше \(1\) или меньше \(-1\).

\subsection{Пример}

Пусть после преобразований получено:

\[
2\sin^2x+2\sin x-1=0.
\]

Делаем замену:

\[
t=\sin x,\qquad -1\le t\le1.
\]

Получаем:

\[
2t^2+2t-1=0.
\]

Дискриминант:

\[
D=2^2-4\cdot2\cdot(-1)=12.
\]

\[
t_{1,2}=\frac{-2\pm\sqrt{12}}{4}
=\frac{-2\pm2\sqrt3}{4}
=\frac{-1\pm\sqrt3}{2}.
\]

Проверяем значения:

\[
t_1=\frac{-1+\sqrt3}{2}\in[-1;1],
\]

\[
t_2=\frac{-1-\sqrt3}{2}<-1.
\]

Значит, второе значение не подходит, а первое даёт:

\[
\sin x=\frac{\sqrt3-1}{2}.
\]

\section{Отбор корней по ограничению}

Если ограничение записано для одной функции, а после решения получена другая функция, нужно отбирать решения на окружности.

Например, если получено:

\[
\sin x=a,\qquad 0<a<1,
\]

а ограничение:

\[
\cos x\le0,
\]

то из двух решений за период нужно взять то, которое лежит во второй четверти.

\[
x=\arcsin a+2\pi n
\]

лежит в первой четверти и не подходит, потому что там \(\cos x>0\).

\[
x=\pi-\arcsin a+2\pi n
\]

лежит во второй четверти и подходит, потому что там \(\cos x<0\).

\begin{center}
\begin{tikzpicture}[scale=1.15,>=Stealth]
  \draw[->] (-1.35,0)--(1.45,0) node[right] {\(\cos x\)};
  \draw[->] (0,-1.35)--(0,1.45) node[above] {\(\sin x\)};
  \draw[darkaccent,line width=0.9pt] (0,0) circle (1);
  \draw[accent,line width=2pt] (0,1) arc (90:270:1);

  \coordinate (A) at ({cos(35)},{sin(35)});
  \coordinate (B) at ({cos(145)},{sin(145)});

  \draw[textgray,dashed] (0,0)--(A);
  \draw[accent,dashed] (0,0)--(B);

  \fill[textgray] (A) circle (1.8pt);
  \fill[accent] (B) circle (2pt);

  \node[above right] at (A) {\(\arcsin a\)};
  \node[above left] at (B) {\(\pi-\arcsin a\)};
  \node[left] at (-1.2,-0.2) {\small \(\cos x\le0\)};
\end{tikzpicture}
\end{center}

\section{Уравнения с отрицательным значением косинуса}

Если получено:

\[
\cos x=-a,\qquad 0<a<1,
\]

то лучше сначала вынести минус отдельно:

\[
\cos x=-a.
\]

Тогда решения можно записать так:

\[
x=\pm\arccos(-a)+2\pi n,\qquad n\in\Z.
\]

Но часто удобнее использовать положительное число \(a\):

\[
\cos x=-a
\quad\Longleftrightarrow\quad
x=\pi\pm\arccos a+2\pi n,\qquad n\in\Z.
\]

Например:

\[
\cos x=\frac{2-\sqrt{26}}{2}.
\]

Это число отрицательно, потому что \(\sqrt{26}>2\). Выносим минус:

\[
\frac{2-\sqrt{26}}{2}
=
-\frac{\sqrt{26}-2}{2}.
\]

Значит:

\[
\cos x=-\frac{\sqrt{26}-2}{2}.
\]

Ответ записываем через положительное число:

\[
x=\pi\pm\arccos\frac{\sqrt{26}-2}{2}+2\pi n,\qquad n\in\Z.
\]

\section{Уравнения с модулем}

\subsection{Определение модуля}

Модуль раскрывается по правилу:

\[
|A|=
\begin{cases}
A, & A\ge0,\\
-A, & A<0.
\end{cases}
\]

Поэтому уравнение с модулем обычно разбивается на две системы.

\subsection{Пример}

Решим уравнение:

\[
|\sin x|=\cos x.
\]

Раскрываем модуль по определению:

\[
\begin{cases}
\sin x\ge0,\\
\sin x=\cos x
\end{cases}
\quad\text{или}\quad
\begin{cases}
\sin x<0,\\
-\sin x=\cos x.
\end{cases}
\]

Первая система:

\[
\sin x=\cos x
\quad\Longleftrightarrow\quad
\tg x=1.
\]

\[
x=\frac{\pi}{4}+\pi n,\qquad n\in\Z.
\]

С учётом \(\sin x\ge0\) подходит:

\[
x=\frac{\pi}{4}+2\pi n.
\]

Вторая система:

\[
-\sin x=\cos x
\quad\Longleftrightarrow\quad
\tg x=-1.
\]

\[
x=-\frac{\pi}{4}+\pi n,\qquad n\in\Z.
\]

С учётом \(\sin x<0\) подходит:

\[
x=\frac{7\pi}{4}+2\pi n.
\]

Ответ:

\[
x=\frac{\pi}{4}+2\pi n
\quad\text{или}\quad
x=\frac{7\pi}{4}+2\pi n,\qquad n\in\Z.
\]

\section{Когда вместо корней лучше использовать понижение степени}

Если в уравнении появляется \(\sin^2x\) или \(\cos^2x\), иногда лучше не извлекать квадратный корень, а использовать формулы понижения степени.

\[
\sin^2x=\frac{1-\cos2x}{2},
\qquad
\cos^2x=\frac{1+\cos2x}{2}.
\]

Например:

\[
\sin^2x=\frac12.
\]

Можно записать:

\[
\frac{1-\cos2x}{2}=\frac12.
\]

Отсюда:

\[
1-\cos2x=1,
\]

\[
\cos2x=0.
\]

Тогда:

\[
2x=\frac{\pi}{2}+\pi n,
\]

\[
x=\frac{\pi}{4}+\frac{\pi n}{2},\qquad n\in\Z.
\]

Это короче, чем отдельно решать

\[
\sin x=\frac1{\sqrt2}
\quad\text{и}\quad
\sin x=-\frac1{\sqrt2}.
\]

\section{Оформление решения на ЕГЭ}

\begin{itemize}
  \item Не пишите сразу только один подходящий корень, если перед этим уравнение даёт два корня. Сначала покажите оба, затем объясните отбор.
  \item Не используйте неясную запись \(\pm\), если она скрывает разные случаи отбора.
  \item Если делите на \(\sin x\) или \(\cos x\), сначала объясните, почему делитель не равен нулю.
  \item Если получаете значение вида \(\frac{\sqrt5-1}{2}\), отдельно проверьте, что оно лежит в промежутке \([-1;1]\).
  \item Если значение отрицательное и записывается через корень, удобно вынести минус и работать с положительным числом.
  \item Пишите решение в столбик: так ясно видно, что из чего следует.
\end{itemize}

\section{Типичные ошибки}

\begin{enumerate}
  \item \textbf{Возвести в квадрат без ограничения.}

  Нельзя переходить от \(\sqrt{f(x)}=g(x)\) сразу к \(f(x)=g^2(x)\). Нужно добавить \(g(x)\ge0\).

  \item \textbf{Поставить ограничение не на ту часть.}

  В уравнении \(\sqrt{f(x)}=g(x)\) основное ограничение ставится на \(g(x)\), а не на \(f(x)\).

  \item \textbf{Забыть диапазон синуса и косинуса.}

  Значения \(\sin x\) и \(\cos x\) лежат только в промежутке \([-1;1]\).

  \item \textbf{Не отобрать корни по ограничению.}

  Если после возведения в квадрат получено тригонометрическое уравнение, его корни нужно проверить по системе.

  \item \textbf{Потерять ветку при раскрытии модуля.}

  Модуль раскрывается двумя случаями, а не одним.

  \item \textbf{Делить на \(\sin x\) или \(\cos x\) без проверки.}

  Деление на ноль недопустимо. Нужно заранее исключить нулевой случай.

  \item \textbf{Неверно записать ответ при отрицательном значении косинуса.}

  Для \(\cos x=-a\), где \(0<a<1\), удобно писать
  \[
  x=\pi\pm\arccos a+2\pi n.
  \]

  \item \textbf{Считать, что найденное число ``очевидно'' подходит.}

  Для выражений с корнями нужно строго сравнивать с \(1\), \(-1\), \(0\) или другим ограничением.
\end{enumerate}

\section{Короткий чек-лист перед ответом}

Перед записью ответа проверьте:

\begin{itemize}
  \item все ограничения записаны до возведения в квадрат;
  \item после замены \(t=\sin x\) или \(t=\cos x\) учтено \(-1\le t\le1\);
  \item все найденные значения проверены по ограничениям;
  \item при модуле рассмотрены оба случая;
  \item если была тригонометрическая окружность, выбран правильный сектор;
  \item если значение отрицательное, ответ записан корректно;
  \item в ответе указано \(n\in\Z\).
\end{itemize}

\newpage

\section*{Тренировочный блок}

Решите уравнения. Ответы запишите в общем виде.

\subsection*{Средний уровень}

\begin{enumerate}
  \item
  \[
  \sqrt{\sin x}=\sqrt{\cos x}.
  \]

  \item
  \[
  \sqrt{\cos^2x}=\cos x.
  \]

  \item
  \[
  |\sin x|=\sin x.
  \]
\end{enumerate}

\subsection*{Выше среднего уровня}

\begin{enumerate}[resume]
  \item
  \[
  \sqrt{\sin x}=-\cos x.
  \]

  \item
  \[
  \sqrt{\cos x}=\sin x.
  \]

  \item
  \[
  \sqrt{\sin^2x}=\frac12.
  \]

  \item
  \[
  |\cos x|=\sin x.
  \]

  \item
  \[
  \sqrt{2\sin x+1}=\cos x.
  \]

  \item
  \[
  \sqrt{\cos^2x}=\frac12-\cos x.
  \]
\end{enumerate}

\subsection*{Сложный уровень}

\begin{enumerate}[resume]
  \item
  \[
  \sqrt{1+\sin x}=|\cos x|.
  \]
\end{enumerate}

\newpage

\section*{Ответы}

\renewcommand{\arraystretch}{1.45}

\begin{table}[h!]
\centering
\caption{Ответы к тренировочному блоку}
\begin{tabularx}{\linewidth}{cY}
\toprule
№ & Ответ \\
\midrule
1 & \ans{x=\frac{\pi}{4}+2\pi n,\quad n\in\Z} \\

2 & \ans{x\in\left[-\frac{\pi}{2}+2\pi n;\frac{\pi}{2}+2\pi n\right],\quad n\in\Z} \\

3 & \ans{x\in[2\pi n;\pi+2\pi n],\quad n\in\Z} \\

4 & \ans{x=\frac{3\pi}{4}+2\pi n,\quad n\in\Z} \\

5 & \ans{x=\frac{\pi}{6}+2\pi n,\quad n\in\Z} \\

6 & \ans{x=\frac{\pi}{6}+\pi n\quad\text{или}\quad x=\frac{5\pi}{6}+\pi n,\quad n\in\Z} \\

7 & \ans{x=\frac{\pi}{4}+2\pi n\quad\text{или}\quad x=\frac{3\pi}{4}+2\pi n,\quad n\in\Z} \\

8 & \ans{x=2\pi n\quad\text{или}\quad x=\pi-\arcsin\frac{\sqrt5-1}{2}+2\pi n,\quad n\in\Z} \\

9 & \ans{x=\pm\arccos\frac14+2\pi n,\quad n\in\Z} \\

10 & \ans{x=2\pi n,\quad n\in\Z} \\
\bottomrule
\end{tabularx}
\end{table}

\end{document}